%! Random Sampling Methods — Unit 3, Lesson 3.2 — Guided Notes & Practice
% THE in-class document, in the MAIN IDEAS / NOTES shape (see LESSON_SHAPE.md §1), at the
% COMPACT budget: the vocabulary box, then ONE guidednotes table of FOUR rows (three instruction
% rows + Guided Practice), 14 problems.
%   I DO   : the teacher fills the definition lines and works the FIRST problem of each grid
%            (1, 4, 7).
%   WE DO  : the class works the rest; the last row, Guided Practice (11–14), is worked together.
%   YOU DO : nothing on this page — ../homework/ IS the individual practice.
% THE TRAP is problem 6 (one grade picked at random: every student has a 1-in-4 chance, but it
% is not an SRS — a group that mixes grades can never be chosen; EK DAT-2.C.2).
% THE CRUX is problem 10, in the last instruction row: "use the grades as clusters" — strata are
% groups alike inside on something related to the response, and we sample from EVERY one;
% clusters each look like the whole population, and we take ALL of SOME (EK DAT-2.C.3, C.4).
% Systematic random sampling gets one line in row 2 and no problems (teacher's decision).
%
% THE CONTEXTS (instruction, guided practice, and homework each use a different one)
%  Rows 1–3 — Lakeview High: 1,200 students, 300 in each of grades 9–12, and 48 advisory groups
%          that each mix all four grades. Question: mean hours per week worked at a paid job
%          (seniors work far more than freshmen, so the grades differ on the response).
%  Row 4 (GP) — Riverside Apartments: 480 apartments in 40 buildings, each building holding
%          4 one-, 4 two-, and 4 three-bedroom apartments. Question: mean cars per household.
%  Homework — a hospital's nurses (shifts and floors). AP Practice — farms, textbooks, an airplane.
%
% PAGE LOCKSTEP with notes_key/: \termblank -> \termans, \blank -> \ans, 4th arg of \pcell.
% TABLE MECHANICS: inside a cell `\\` ENDS THE ROW — use \par. Long rows are split into
% sub-rows (a bare `\\` then `& ...`) so the table can break between them.
% NO SKETCH QUESTIONS: every display is given, read and used.
\documentclass[12pt]{article}
\usepackage{apstats-article}
\usepackage{apstats-boxes}

\begin{document}

\pageheader{Unit 3, Lesson 3.2}{Guided Notes \& Practice}
% NAMESTRIP: no \namedateperiod here — the name row lives on the cover only.

\vspace{2pt}

\boxguard[16]
\begin{vocabbox}
\small
Fill in each term \textbf{as we name it} in the notes below.
\par
\termblank{Simple random sample}
\par
\termblank{Without replacement}
\par
\termblank{Stratified random sample}
\par
\termblank{Cluster sample}
\par
\end{vocabbox}

\small
\begin{guidednotes}
% ── ROW 1: the simple random sample ──────────────────────────────────────────
\mainidea[Chance chooses]{SRS} &
In a \textbf{simple random sample (SRS)} of size $n$, every \blank{1.4cm} of $n$ individuals
has an \blank{1.3cm} chance of being the sample. Sampling \textbf{without replacement}: an
individual can be chosen only \blank{1.2cm}. \textbf{With replacement}: an individual can be
chosen more than once.

\vspace{3pt}
\stepnum{1}\ \textbf{Label:} give each individual a different number, 1 to $N$.\par
\stepnum{2}\ \textbf{Randomize:} use a random number generator or a table of random digits.\par
\stepnum{3}\ \textbf{Select:} skip numbers out of range and \blank{1.6cm} until you have $n$.
\\
& Lakeview High has 1{,}200 students: 300 in each of grades 9--12. The principal wants the mean
number of hours per week Lakeview students work at a paid job.
\notesprompt{Choose an SRS of Lakeview students.}
\begin{probgrid}{|Y|Y|Y|}
\pcell{1}{Describe how to use a random number generator to choose an SRS of 80 students.}{2.8cm}{} &
\pcell{2}{Students are numbered 0001--1200. Use these digits, four at a time, to pick the first three students:\par
\texttt{0738 1195 9661 0738 0452}}{2.8cm}{} &
\pcell{3}{A teacher draws a name from a hat of all 1{,}200, puts it back, then draws again. With or without replacement?}{2.8cm}{} \\ \hline
\end{probgrid}
\\ \hline
% ── ROW 2: strata and clusters — THE TRAP ────────────────────────────────────
\mainidea[Groups first]{Strata \& clusters} &
\textbf{Stratified random sample:} split the population into \textbf{strata}, groups whose
members are \blank{1.2cm} each other; take an \blank{1.0cm} from \blank{1.1cm} stratum and
combine them.\par\vspace{3pt}
\textbf{Cluster sample:} split the population into \textbf{clusters}, groups that each look like
a small copy of the \blank{2.0cm}; randomly choose some clusters and measure \blank{1.2cm}
individual in them.\par\vspace{3pt}
{\footnotesize (One more, named only: a \textbf{systematic random sample} picks a random starting
point on a list, then every $k$th individual.)}
\\
& \notesprompt{Lakeview's 48 advisory groups each mix students from all four grades. Name each method.}
\begin{probgrid}{|Y|Y|Y|}
\pcell{4}{Randomly choose 20 students from each of grades 9, 10, 11, and 12.}{1.4cm}{} &
\pcell{5}{Randomly choose 2 of the 48 advisory groups; survey every student in them.}{1.4cm}{} &
\pcell{6}{Randomly choose 1 grade; survey all 300 in it. Each student has a 1-in-4 chance. Is it an SRS?}{1.4cm}{} \\ \hline
\end{probgrid}
\\ \hline
% ── ROW 3: choosing a method — THE CRUX ──────────────────────────────────────
\mainidea[Which plan?]{Choosing} &
\textbf{Stratify} when the response \blank{1.4cm} from group to group: strata alike inside give
an estimate that \blank{1.6cm} less from sample to sample than an SRS of the same size.
\textbf{Cluster} when reaching scattered individuals is \blank{2.0cm}: it is cheaper and more
practical, and it works best when each cluster looks like the population.
\notesprompt{Seniors at Lakeview work far more hours than freshmen. Choose and defend a plan.}
\begin{probgrid}{|Y|Y|}
\pcell{7}{Why would 20 from each grade beat an SRS of 80 for estimating mean hours worked?}{1.6cm}{} &
\pcell{8}{Give one advantage of the advisory-group cluster sample over the plan in problem~7.}{1.6cm}{} \\ \hline
\pcell{9}{Should the principal use the 48 advisory groups as \emph{strata}? Why or why not?}{1.9cm}{} &
\pcell{10}{Jalen: ``Use the grades as clusters --- randomly pick 2 grades and survey everyone in them.'' What is wrong with his plan?}{1.9cm}{} \\ \hline
\end{probgrid}
\\
& \fcolorbox{navy}{sky}{\parbox{\dimexpr\linewidth-2\fboxsep-2\fboxrule\relax}{\centering
\bfseries Strata are alike inside: sample from \blank{1.2cm} one.\par
Clusters each look like the whole: take \blank{0.8cm} of \blank{1.1cm} of them.}}
\\ \hline
% ── ROW 4: guided practice — WE DO, together, fresh context ──────────────────
\mainidea[Guided practice]{Parking at Riverside} &
Riverside Apartments has 480 apartments in 40 buildings. The manager wants the mean number of
cars per household to plan new parking.
\par\vspace{3pt}
\begin{minipage}{\linewidth}\centering\footnotesize
\setlength{\tabcolsep}{4pt}%
\renewcommand{\arraystretch}{1.2}
\begin{tabular}{@{} l c c @{}}
\rowcolor{sky}
\textbf{Apartment type} & \textbf{In each building} & \textbf{In the complex} \\
\hline
One-bedroom   & 4 & 160 \\
Two-bedroom   & 4 & 160 \\
Three-bedroom & 4 & 160 \\
\end{tabular}
\end{minipage}
\\
& \notesprompt{We work these together.}
\begin{probgrid}{|Y|Y|}
\pcell{11}{Describe how to choose an SRS of 60 apartments with a random number generator.}{2.2cm}{} &
\pcell{12}{Plan B: randomly choose 5 buildings and survey every apartment in them. Name it; give one advantage.}{2.2cm}{} \\ \hline
\pcell{13}{Plan C: an SRS of 20 apartments of each type. Name it. Why might it beat an SRS of 60?}{2.2cm}{} &
\pcell{14}{Plan D: ``Treat the three types as clusters: pick one type at random and survey all 160.'' Respond.}{2.2cm}{} \\ \hline
\end{probgrid}
\\ \hline
\end{guidednotes}

% ── THE NOTES END AT GUIDED PRACTICE ─────────────────────────────────────────
% There is deliberately no "you do" here: ../homework/ is the individual practice,
% started in class in the last 8 minutes while the teacher circulates.

\end{document}
